% ============================================================================
\section{Testing}
\label{sec:testing}
% ============================================================================

\subsection{Unit tests of the decision logic}
\file{firmware/test/host/test_core_logic.c} compiles the pure \code{core_logic} and the type
helpers with the PC's gcc (or a gcc Docker image) using \code{-Wall -Wextra -Werror}, and checks
them against expected results. Run it with \code{make test}.

\begin{center}
\small
\begin{tabularx}{\linewidth}{@{}lL@{}}
\toprule
Group & What it proves \\
\midrule
Schedule & Opening and closing edges, closed weekdays, overnight windows, minutes to close,
  invalid schedules rejected. \\
Access policy & Every rule in order: unknown, disabled, anti-passback both ways, lockdown,
  in hours, privileged after hours, escort exit, after-hours refusal. \\
Room logic & Occupancy and grace period, dim warning, daylight thresholds and hysteresis, manual
  ON/OFF override, energy used/saved, the daylight change flag. \\
Admin login & Happy path, non-admin card, wrong PIN, lockout after 3 failures and its expiry,
  step timeouts, late answers, session expiry, millisecond-counter wrap-around. \\
Types & Parsing of role names and card UIDs, including invalid input. \\
\bottomrule
\end{tabularx}
\end{center}
\textbf{Result: 70 checks, 0 failures.}

\subsection{Build}
\code{idf.py build} for the ESP32-S3 with ESP-IDF~v6.0.3 completes with \textbf{0 warnings}. The
application image is 0x628C0 bytes (about 395\,KB), leaving 74\,\% of the application partition
free.

\subsection{Simulation testing in Wokwi}
The system was run in the Wokwi simulator through VS Code. Observed results:

\begin{center}
\small
\begin{tabularx}{\linewidth}{@{}lL@{}}
\toprule
Scenario & Result \\
\midrule
Boot & RTC time read, SD card mounted, SHA-256 self-test passed, 4 demo cards seeded, both
  rooms reported online. \\
Card scans & Alice, Bob and Carol entered and were logged as \code{ACCESS ... ok=1 reason=OK}. \\
Reboot & The log resumed after the last sequence number; the ledger reported
  \enquote{2 people inside at boot}. \\
Admin login & Card then PIN gave \code{AUTH result=CARD_OK} followed by \code{result=OK}. \\
Light sensor & Moving the slider below 150\,lux logged \code{ROOM ... daylight=0}; above 400\,lux
  \code{daylight=1}. \\
Energy & \code{ENERGY} lines appeared every 10 minutes for both rooms. \\
\bottomrule
\end{tabularx}
\end{center}

\subsection{Defects found and fixed}
Testing in the simulator and code review found real bugs, which were fixed:
\begin{itemize}
  \item \textbf{Crash loop when a lamp faded on.} The ESP32's hardware LED-fade interrupt fired
        endlessly in Wokwi and tripped the interrupt watchdog. Because a software reset does not
        clear that peripheral, every reboot crashed again. The backtrace pointed to
        \code{ledc_fade_isr}; we removed the hardware fade and fade the lamp in software from the
        room task instead.
  \item \textbf{Bus ordering race} (\cref{sec:sync}): fixed with a per-sender mutex.
  \item \textbf{PIN entry}: the PIN was made exactly 4 digits with automatic submit, and a keypad
        repeat guard stopped a single press being counted twice.
  \item \textbf{LCD overflow}: every row is now formatted to exactly 20 characters.
  \item \textbf{Late PIN after a timeout} was accepted by mistake; the state machine now checks
        expiry first, and a unit test covers it.
\end{itemize}

\subsection{Dashboard}
\code{prisma validate}, the TypeScript check and \code{next build} pass. Against a local
PostgreSQL 16 database, the sample sender inserted 18 events, then 18 duplicates on a second run
(idempotent). A bad signature and a stale timestamp were rejected with HTTP 401, and a malformed
body with HTTP 400.

\todo{after the final demo run, add any further scenarios tested (schedule change, after-hours
exit and escort, break-glass, node offline, log verify)}
